\chapter{Expansion as the primary limit tool}

\begin{orient}
\textbf{What this chapter is for.} A Taylor expansion converts a transcendental function into a polynomial plus a controlled remainder, and once every ingredient of a limit is a polynomial the limit is arithmetic. That is why expansion is not one technique among several: it is the general-purpose engine, and almost every other method at a point is either a special case of it or a slower route to the same coefficient. This chapter installs the standard Maclaurin table, teaches order bookkeeping so that you decide in advance how many terms each factor needs rather than discovering it after three failed attempts, and shows how to move the expansion point to a finite $a$ or to infinity. It closes with L'Hopital's rule treated honestly: what its hypotheses actually require, where repeated application is genuinely the fastest route, where it cycles, and the standard limits at which the rule reports failure while the limit exists perfectly well. At the end you should reach for expansion first, use L'Hopital deliberately rather than reflexively, and never be surprised by how far you had to expand.
\end{orient}

\section{The engine, and the table it runs on}

The mechanism is one line. If $f$ is $n$ times differentiable at $0$, then
\[
f(x)=\sum_{k=0}^{n}\frac{f^{(k)}(0)}{k!}x^{k}+o(x^{n})\qquad(x\to0),
\]
and if $f$ is $(n+1)$ times differentiable near $0$ the remainder can be written in Lagrange form as $\frac{f^{(n+1)}(\xi)}{(n+1)!}x^{n+1}$ for some $\xi$ between $0$ and $x$. For limit work the little-o form is the one you need, because the only question ever asked is whether the discarded tail dies faster than the denominator.

The reason this dominates the subject is structural. A limit at a point is a question about the first few coefficients of the functions involved, and Taylor's theorem is the statement that those coefficients exist and are computable. Every competing technique extracts the same coefficients less directly. L'Hopital's rule computes $f^{(k)}(0)$ and $g^{(k)}(0)$ one at a time by differentiating the whole expression, which reintroduces the product and quotient rules at every step. Asymptotic equivalence extracts only the first coefficient and then has nothing further to say. Expansion extracts all of them at once, from a table you already know, and the bookkeeping is addition of polynomials.

\begin{keynote}[L'Hopital is a corollary, not a rival]
Suppose $f$ and $g$ vanish at $a$ and are twice differentiable there. Expanding both,
\[
\frac{f(x)}{g(x)}=\frac{f'(a)(x-a)+\tfrac12f''(a)(x-a)^{2}+o((x-a)^{2})}{g'(a)(x-a)+\tfrac12g''(a)(x-a)^{2}+o((x-a)^{2})}
=\frac{f'(a)+\tfrac12f''(a)(x-a)+o(x-a)}{g'(a)+\tfrac12g''(a)(x-a)+o(x-a)},
\]
and if $g'(a)\neq0$ the limit is $f'(a)/g'(a)$. That is exactly L'Hopital's conclusion, obtained by cancelling one power of $(x-a)$ from a pair of expansions. One round of the rule is therefore the same operation as truncating both series at order $1$; $k$ rounds is truncating at order $k$. The difference is that expansion performs all $k$ truncations in a single pass, from a table, whereas the rule performs them one at a time, by differentiating expressions that grow at every step. Everything in the last section of this chapter follows from that observation.
\end{keynote}

You cannot use the engine without the fuel. The following table is the minimum, given to the order that actually comes up. The final column is the leading deviation from the linear or constant part, which is the number that decides most problems and the number most often misremembered.

\begin{center}
\footnotesize
\setlength{\tabcolsep}{5pt}
\renewcommand{\arraystretch}{1.32}
\begin{tabular}{@{}lll@{}}
\toprule
\mh{Function} & \mh{Maclaurin series, to the order that matters} & \mh{Signature term} \\
\midrule
$\eu^{x}$ & $1+x+\tfrac{x^{2}}{2}+\tfrac{x^{3}}{6}+\tfrac{x^{4}}{24}+O(x^{5})$ & $+\tfrac{x^{2}}{2}$ \\
$\ln(1+x)$ & $x-\tfrac{x^{2}}{2}+\tfrac{x^{3}}{3}-\tfrac{x^{4}}{4}+O(x^{5})$ & $-\tfrac{x^{2}}{2}$ \\
$(1+x)^{a}$ & $1+ax+\tfrac{a(a-1)}{2}x^{2}+\tfrac{a(a-1)(a-2)}{6}x^{3}+O(x^{4})$ & $+\tfrac{a(a-1)}{2}x^{2}$ \\
$\sin x$ & $x-\tfrac{x^{3}}{6}+\tfrac{x^{5}}{120}+O(x^{7})$ & $-\tfrac{x^{3}}{6}$ \\
$\cos x$ & $1-\tfrac{x^{2}}{2}+\tfrac{x^{4}}{24}-\tfrac{x^{6}}{720}+O(x^{8})$ & $-\tfrac{x^{2}}{2}$ \\
$\tan x$ & $x+\tfrac{x^{3}}{3}+\tfrac{2x^{5}}{15}+\tfrac{17x^{7}}{315}+O(x^{9})$ & $+\tfrac{x^{3}}{3}$ \\
$\sec x$ & $1+\tfrac{x^{2}}{2}+\tfrac{5x^{4}}{24}+\tfrac{61x^{6}}{720}+O(x^{8})$ & $+\tfrac{x^{2}}{2}$ \\
$\arcsin x$ & $x+\tfrac{x^{3}}{6}+\tfrac{3x^{5}}{40}+\tfrac{5x^{7}}{112}+O(x^{9})$ & $+\tfrac{x^{3}}{6}$ \\
$\arctan x$ & $x-\tfrac{x^{3}}{3}+\tfrac{x^{5}}{5}-\tfrac{x^{7}}{7}+O(x^{9})$ & $-\tfrac{x^{3}}{3}$ \\
$\sinh x$ & $x+\tfrac{x^{3}}{6}+\tfrac{x^{5}}{120}+O(x^{7})$ & $+\tfrac{x^{3}}{6}$ \\
$\cosh x$ & $1+\tfrac{x^{2}}{2}+\tfrac{x^{4}}{24}+\tfrac{x^{6}}{720}+O(x^{8})$ & $+\tfrac{x^{2}}{2}$ \\
$\tanh x$ & $x-\tfrac{x^{3}}{3}+\tfrac{2x^{5}}{15}-\tfrac{17x^{7}}{315}+O(x^{9})$ & $-\tfrac{x^{3}}{3}$ \\
$\operatorname{arsinh}x$ & $x-\tfrac{x^{3}}{6}+\tfrac{3x^{5}}{40}+O(x^{7})$ & $-\tfrac{x^{3}}{6}$ \\
\bottomrule
\end{tabular}
\end{center}

Three structural facts make this table smaller than it looks. The hyperbolic entries are the circular ones with the alternating signs removed, so $\sinh$, $\cosh$ and $\tanh$ cost nothing extra once $\sin$, $\cos$ and $\tan$ are known. The signature term of a function and of its inverse always have opposite signs and, for the sine family, equal magnitude: $\sin$ has $-\tfrac{x^{3}}{6}$ and $\arcsin$ has $+\tfrac{x^{3}}{6}$, $\tan$ has $+\tfrac{x^{3}}{3}$ and $\arctan$ has $-\tfrac{x^{3}}{3}$, $\sinh$ has $+\tfrac{x^{3}}{6}$ and $\operatorname{arsinh}$ has $-\tfrac{x^{3}}{6}$. And $(1+x)^{a}$ subsumes every radical, so no separate root entries are needed.

\tech{The standard expansion table}{easy}
\cue Any limit at all as $x\to0$, and any limit that a substitution can move to $0$.
\method Know the table above cold, to the orders shown. In use, replace each transcendental factor by its polynomial and carry an explicit remainder. The single most valuable line of the table is the group of four functions that all begin $x\pm\tfrac{x^{3}}{k}$, since differences within that group are the commonest engineered limits in the subject.

\begin{worked}
\[
\lim_{x\to0}\frac{\arcsin x-\arctan x}{x^{3}}.
\]
The linear terms cancel by construction, and the signature terms are $+\tfrac{x^{3}}{6}$ and $-\tfrac{x^{3}}{3}$. So the numerator is $\bigl(\tfrac16+\tfrac13\bigr)x^{3}+O(x^{5})=\tfrac{x^{3}}{2}+O(x^{5})$ and the limit is $\dfrac12$, verified as $0.4999875$ at $x=10^{-2}$. The whole calculation is two table lookups and one addition; L'Hopital would need three rounds, and the third derivative of $\arcsin$ is not a pleasant object.

The same two lookups settle the most-asked limit in the family:
\[
\lim_{x\to0}\frac{\tan x-\sin x}{x^{3}}=\Bigl(\frac13-\bigl(-\frac16\bigr)\Bigr)=\frac12 ,
\]
verified as $0.5000001$ at $x=10^{-3}$. Note what this says about the trap in the previous chapter: $\tan x\sim x$ and $\sin x\sim x$ are both true, and their difference is nonetheless of order $3$ with coefficient $\tfrac12$. Equivalence could not have found this; the table found it in one line.
\end{worked}

\begin{trap}
Confusing the signs within the $x\pm\tfrac{x^{3}}{k}$ family. The reliable rule is the inverse-function rule above: $\tan$ and $\arcsin$ bend up, $\arctan$ and $\tanh$ bend down, $\sin$ is negative and $\sinh$ is positive. The second trap is quoting $\arcsin x=x+\tfrac{x^{3}}{3}$, borrowing $\tan$'s denominator; it is $6$, not $3$.
\end{trap}

\begin{seealsob}Expand to the first surviving order; the binomial series for radicals; composite and nested expansion.\end{seealsob}

The table covers functions of $x$. Two extensions cover essentially everything else: an arbitrary real power, which absorbs all radicals, and a function of a function, which absorbs all the composites. Take the radicals first, because they are the ones most often mishandled by stopping one term too early.

\tech{The binomial series for radicals}{easy}
\cue A square root, cube root, or any fractional power of something near $1$; equivalently, any radical whose argument can be written as $c\,(1+u)$ with $u\to0$.
\method Use
\[
(1+u)^{a}=1+au+\frac{a(a-1)}{2}u^{2}+\frac{a(a-1)(a-2)}{6}u^{3}+O(u^{4}),
\]
where $u$ is the \emph{whole} small perturbation and may itself be a series. Force the argument into the form $1+u$ first, pulling out constants and dominant powers, since the series is valid only there.

\begin{worked}
\[
\lim_{x\to0}\frac{\sqrt{1+2x}-\sqrt[3]{1+3x}}{x^{2}}.
\]
With $a=\tfrac12$, $u=2x$: $\sqrt{1+2x}=1+x-\tfrac{x^{2}}{2}+O(x^{3})$, since $\tfrac{a(a-1)}{2}u^{2}=\tfrac{(1/2)(-1/2)}{2}\cdot4x^{2}=-\tfrac{x^{2}}{2}$. With $a=\tfrac13$, $u=3x$: $\sqrt[3]{1+3x}=1+x-x^{2}+O(x^{3})$, since $\tfrac{(1/3)(-2/3)}{2}\cdot9x^{2}=-x^{2}$. The linear terms cancel and the limit is $-\tfrac12-(-1)=\dfrac12$. Verified as $0.49884$ at $x=10^{-3}$.

A version at infinity, where the perturbation has to be manufactured:
\[
\lim_{x\to\infty}x^{3}\Bigl(\sqrt[3]{x^{3}+3x}-\sqrt{x^{2}+2}\Bigr)
=\lim_{x\to\infty}x^{3}\Bigl[\bigl(x+\tfrac1x-\tfrac1{x^{3}}\bigr)-\bigl(x+\tfrac1x-\tfrac1{2x^{3}}\bigr)\Bigr]=-\frac12 ,
\]
both radicals having been written as $x(1+u)^{a}$ with $u=3x^{-2}$ and $u=2x^{-2}$. Verified to six figures at $x=10^{3}$ with $50$-digit arithmetic; ordinary floating point cannot see it, because the two radicals agree to twelve significant figures and the difference is pure round-off.
\end{worked}

\begin{trap}
Using $\sqrt{1+u}\approx1+\tfrac{u}{2}$ when the problem needs the $-\tfrac{u^{2}}{8}$ term. A radical limit whose denominator is $x^{2}$ always needs second order, and a limit like the one at infinity above needs third. The other error is applying the series to $\sqrt{x^{2}+2}$ directly, as though $x^{2}$ were small; the perturbation must be the small quantity, which is why the $x$ is pulled out first.
\end{trap}

\begin{seealsob}The standard expansion table; composite and nested expansion; expansion at infinity via $t=1/x$.\end{seealsob}

Radicals were the first composition. The general case is a series inside a series, and it looks more frightening than it is, because the inner function vanishes at the origin and that fact alone bounds the work.

\tech{Composite and nested expansion}{medium}
\cue A series inside a series: $\ln(\cos x)$, $\eu^{\sin x}$, $\sin(\tan x)$, $\sqrt{1+\sin^{2}x}$, $\tan(\ln(1+x))$.
\method Expand the inner function first, call the result $u=u(x)$ with $u=O(x)$, then substitute into the outer expansion in powers of $u$ and discard every term whose total degree in $x$ exceeds the target. The work is bounded because the $u^{k}$ term of the outer series contributes only from degree $k$ upward: if the target is degree $N$ and $u$ has order $1$, only the first $N$ terms of the outer series can matter, and usually far fewer.

\begin{worked}
\[
\lim_{x\to0}\frac{\sqrt[3]{1+\sin^{2}x}-\sqrt{\cos x}}{x^{2}}.
\]
Target degree $2$. Inner functions: $\sin^{2}x=x^{2}+O(x^{4})$ and $\cos x-1=-\tfrac{x^{2}}{2}+O(x^{4})$, both of order $2$, so only the linear term of each outer binomial can contribute. Hence $\sqrt[3]{1+\sin^{2}x}=1+\tfrac13x^{2}+O(x^{4})$ and $\sqrt{\cos x}=1-\tfrac14x^{2}+O(x^{4})$, and the numerator is $\bigl(\tfrac13+\tfrac14\bigr)x^{2}+O(x^{4})$. The limit is $\dfrac{7}{12}$, verified as $0.5833331$ at $x=10^{-3}$.

A composite where the outer function must be exploited rather than expanded:
\[
\lim_{x\to0}\frac{\eu^{\sin x}-\eu^{\tan x}}{x^{3}}
=\lim_{x\to0}\eu^{\sin x}\cdot\frac{1-\eu^{\tan x-\sin x}}{x^{3}}
=\lim_{x\to0}\frac{-(\tan x-\sin x)}{x^{3}}=-\frac12 ,
\]
since $\tan x-\sin x=\tfrac{x^{3}}{2}+O(x^{5})$ and $\eu^{u}-1\sim u$. Verified as $-0.50504$ at $x=10^{-2}$, converging to $-0.5$. Factoring out the common exponential turned a two-series problem into a one-series problem.

A nesting where the outer series must be carried to $u^{2}$:
\[
\lim_{x\to0}\frac{\ln(\cos x)+\tfrac{x^{2}}{2}}{x^{4}}=-\frac1{12}.
\]
Here $u=\cos x-1=-\tfrac{x^{2}}{2}+\tfrac{x^{4}}{24}+O(x^{6})$ has order $2$ and the target degree is $4$, so the outer series $\ln(1+u)=u-\tfrac{u^{2}}{2}+O(u^{3})$ must be taken to $u^{2}$ and no further, since $u^{3}=O(x^{6})$. Then $\ln(\cos x)=-\tfrac{x^{2}}{2}+\tfrac{x^{4}}{24}-\tfrac12\cdot\tfrac{x^{4}}{4}+O(x^{6})=-\tfrac{x^{2}}{2}-\tfrac{x^{4}}{12}+O(x^{6})$. Verified as $-0.0833356$ at $x=10^{-2}$ against $-1/12=-0.0833333$.
\end{worked}

\begin{trap}
Substituting the entire inner series into the entire outer series and multiplying out. That generates dozens of irrelevant high-degree terms and buries the two that matter. Decide the target degree first, then note the order of $u$, then expand the outer series only as far as $u^{\lfloor N/\ord u\rfloor}$.
\end{trap}

\begin{seealsob}The binomial series for radicals; order bookkeeping across a subtraction; when the third order is genuinely needed.\end{seealsob}

\subsection{The same limit four ways}

It is worth doing one limit by every available method, because the comparison is the argument of this chapter in miniature. Take $\displaystyle L=\lim_{x\to0}\frac{\tan x-\sin x}{x^{3}}$.

\emph{By the table.} The signature terms are $+\tfrac{x^{3}}{3}$ and $-\tfrac{x^{3}}{6}$, so the numerator is $\tfrac{x^{3}}{2}+O(x^{5})$ and $L=\tfrac12$. Two lookups.

\emph{By factoring, then the table.} $\tan x-\sin x=\sin x\bigl(\sec x-1\bigr)=\dfrac{\sin x\,(1-\cos x)}{\cos x}$, and now the three factors have known leading behaviour $x$, $\tfrac{x^{2}}{2}$ and $1$, so the numerator is $\tfrac{x^{3}}{2}+O(x^{5})$ again. This route is the one to prefer when the table has been forgotten, since $1-\cos x\sim\tfrac{x^{2}}{2}$ is the only fact used.

\emph{By L'Hopital.} Three rounds. First $\dfrac{\sec^{2}x-\cos x}{3x^{2}}$, still $\tfrac00$; then $\dfrac{2\sec^{2}x\tan x+\sin x}{6x}$, still $\tfrac00$; then $\dfrac{2\sec^{4}x+4\sec^{2}x\tan^{2}x+\cos x}{6}\to\dfrac{2+0+1}{6}=\dfrac12$. Correct, and three chances to lose a factor of two.

\emph{By asymptotic equivalence.} $\tan x\sim x$ and $\sin x\sim x$, so the numerator is ``$\sim x-x=0$'' and $L=0$. This is wrong. Equivalence is not closed under subtraction, and the failure is silent: nothing in the calculation signals that an illegal step was taken.

Four methods, three correct, one fast, and one that fails without warning. That ranking is stable across the whole subject, and it is why the table comes first.

\section{Order bookkeeping: deciding in advance how far to expand}

Beginners expand, discover the answer is $0/0$ again, expand further, and repeat. That loop is avoidable, and avoiding it is most of what separates a fluent solver from a slow one. The rule is arithmetic on orders, where the \emph{order} of a function is the degree of its lowest nonvanishing Maclaurin term: $\ord(\sin x)=1$, $\ord(1-\cos x)=2$, $\ord(\tan x-x)=3$. Three facts govern that arithmetic. Orders add under multiplication, $\ord(fg)=\ord f+\ord g$, with no exceptions. Orders multiply under composition, $\ord(f\circ g)=\ord f\cdot\ord g$ when $g(0)=0$, which is why $\ord\bigl(\sin(x^{2})\bigr)=2$ and $\ord\bigl(1-\cos(x^{2})\bigr)=4$. And under addition $\ord(f\pm g)\geq\min(\ord f,\ord g)$, with equality unless the leading terms cancel, which is the one place where the arithmetic cannot be done by inspection and where every difficult problem in this section lives.

\begin{keynote}[How many terms, decided in advance]
Let $N/D$ be the quotient and let $d=\ord D$. Expand $N$ through degree $d$ inclusive, carrying $O(x^{d+1})$. If $\ord N<d$ the limit is infinite; if $\ord N=d$ the limit is the ratio of the two leading coefficients; if $N$ vanishes through degree $d$ the limit is $0$. For a numerator that is a product $N=\prod_i f_i$ with $\ord f_i=m_i$ and total order $m=\sum_i m_i$, factor $f_i$ must be expanded only through degree
\[
d-(m-m_i),
\]
because the other factors contribute at least $m-m_i$ degrees for free. That formula is the whole of order bookkeeping.
\end{keynote}

\tech{Expand to the first surviving order, and no further}{easy}
\cue A quotient $N(x)/x^{k}$ or $N(x)/D(x)$ as $x\to0$, where $N$ is a product or a sum of standard functions.
\method Compute $d=\ord D$, apply the budget formula above to each factor of $N$, expand exactly that far, and stop. The identity underneath is that a function of order $m$ satisfies $f(x)=\frac{f^{(m)}(0)}{m!}x^{m}+O(x^{m+1})$, so the answer is a single coefficient and everything beyond it is waste.

\begin{worked}
\[
\lim_{x\to0}\frac{(\eu^{x}-1)^{2}\bigl(\ln(1+x)-\sin x\bigr)^{2}}{x^{6}}.
\]
Here $d=6$. The factors are $f_1=(\eu^{x}-1)^{2}$ with $m_1=2$ and $f_2=\bigl(\ln(1+x)-\sin x\bigr)^{2}$ with $m_2=4$, so $m=6$ and the budgets are $6-(6-2)=2$ for $f_1$ and $6-(6-4)=4$ for $f_2$. Within budget: $\eu^{x}-1=x+O(x^{2})$ gives $f_1=x^{2}+O(x^{3})$, and $\ln(1+x)-\sin x=\bigl(x-\tfrac{x^{2}}{2}+\tfrac{x^{3}}{3}\bigr)-\bigl(x-\tfrac{x^{3}}{6}\bigr)+O(x^{4})=-\tfrac{x^{2}}{2}+O(x^{3})$ gives $f_2=\tfrac{x^{4}}{4}+O(x^{5})$. The product is $\tfrac{x^{6}}{4}+O(x^{7})$ and the limit is $\dfrac14$, verified as $0.2475$ at $x=10^{-2}$.

A quotient where the denominator is not a bare power, so $d$ must be computed:
\[
\lim_{x\to0}\frac{(x-\sin x)(\eu^{x}-1)}{x^{2}(1-\cos x)}.
\]
The denominator has order $2+2=4$, so $d=4$. The numerator's factors have orders $3$ and $1$, total $4$, so each needs only its leading term: $x-\sin x=\tfrac{x^{3}}{6}+O(x^{5})$ and $\eu^{x}-1=x+O(x^{2})$, giving $\tfrac{x^{4}}{6}+O(x^{5})$. The denominator's leading coefficient is $\tfrac12$, so the limit is $\tfrac{1/6}{1/2}=\dfrac13$, verified as $0.333500$ at $x=10^{-3}$.
\end{worked}

\begin{trap}
Expanding one factor to order $5$ and another to order $2$ by instinct. The budget formula exists precisely so that this is a computation rather than a guess. The reverse error is equally costly: expanding every factor to order $6$ here would mean carrying the fifth Maclaurin coefficient of $\ln(1+x)$, which cannot affect the answer. A third failure is applying the budget formula with a stale $m_i$: if a factor turns out to have higher order than you assumed, because its leading terms cancelled, every other factor's budget shrinks and must be recomputed.
\end{trap}

\begin{seealsob}The standard expansion table; order bookkeeping across a subtraction; series division and reciprocal series.\end{seealsob}

The budget formula assumes you know the order of each factor. When the numerator is a difference of two functions with the same leading behaviour, you do not: the order of the difference is exactly what the cancellation determines, and getting it wrong is the standard way to produce a confident wrong answer.

\tech{Order bookkeeping across a subtraction}{medium}
\cue Two expressions with identical leading terms are being subtracted, and the answer lives entirely in what survives.
\method Decide in advance how many orders will cancel, then expand both sides one order past that. If the denominator is $x^{k}$, expand everything through $x^{k}$ inclusive with an explicit $O(x^{k+1})$, and then confirm that the surviving coefficient is genuinely nonzero before reporting it. When two functions agree to order $r$, their difference has order at least $r+1$, and the denominator of the problem tells you what order the setter intends.

\begin{worked}
\[
\lim_{x\to0}\frac{\bigl(W_0(x)-x+x^{2}\bigr)(\sin x-x)}{x^{6}},
\]
where $W_0$ is the principal Lambert function, $W_0(x)=x-x^{2}+\tfrac32x^{3}-\tfrac83x^{4}+O(x^{5})$. The first factor is engineered to cancel two orders, leaving $\tfrac32x^{3}+O(x^{4})$, so $m_1=3$. The second is $\sin x-x=-\tfrac{x^{3}}{6}+O(x^{5})$, so $m_2=3$ and $m=6=d$: each factor needs only its leading term. The product is $\tfrac32\cdot\bigl(-\tfrac16\bigr)x^{6}+O(x^{7})$ and the limit is $-\dfrac14$. Verified as $-0.24564$ at $x=10^{-2}$ and $-0.24867$ at $x=3\times10^{-3}$, converging to $-0.25$ from above as the $O(x^{7})$ correction dies.

A shorter instance of the same discipline, with no exotic functions:
\[
\lim_{x\to0}\frac{\sqrt{1+x}-\eu^{x/2}}{x^{2}}=-\frac14 .
\]
Both functions equal $1+\tfrac{x}{2}+\cdots$, so one order cancels and the denominator says to expand through order $2$. From the table, $\sqrt{1+x}=1+\tfrac{x}{2}-\tfrac{x^{2}}{8}+O(x^{3})$ and $\eu^{x/2}=1+\tfrac{x}{2}+\tfrac{x^{2}}{8}+O(x^{3})$: the quadratic coefficients differ only in sign, which is the whole design. The difference is $-\tfrac{x^{2}}{4}+O(x^{3})$ and the limit is $-\tfrac14$, verified as $-0.2499584$ at $x=10^{-3}$.
\end{worked}

\begin{trap}
The cancellation runs deeper than expected and the term you kept was itself cancelled, so the reported answer is a spurious $0$ or $\infty$. Always check that the surviving coefficient is nonzero as a number, not as a symbol: if it is a difference of two fractions, evaluate it. In the example above, a solver who wrote $W_0(x)=x-x^{2}+O(x^{3})$ would have concluded that the first factor is $O(x^{3})$ with an unknown coefficient, which is true but useless.
\end{trap}

\begin{seealsob}Expand to the first surviving order; engineered-cancellation limits; when the third order is genuinely needed.\end{seealsob}

Some problems make the cancellation the entire point, and once you see the design the answer can be read off without any expansion at all.

\tech{Engineered-cancellation limits}{medium}
\cue The numerator visibly subtracts off the first few Maclaurin terms of a standard function, or two standard functions are averaged so that their low-order terms annihilate.
\method Recognise the design. If the numerator is $f(x)$ minus the first $m$ terms of $f$, and the denominator is $x^{m+1}$ or $x^{m+2}$, the answer is simply the next Taylor coefficient of $f$, that is $f^{(m+1)}(0)/(m+1)!$ or $f^{(m+2)}(0)/(m+2)!$ as the case may be, with $0$ if the intervening coefficient vanishes.

\begin{worked}
The bare form first:
\[
\lim_{x\to0}\frac{1-\cos x-\tfrac12x^{2}}{x^{4}}=-\frac1{24},
\]
because the numerator is exactly $-\tfrac{x^{4}}{24}+O(x^{6})$: the constant and quadratic terms of $\cos$ have been removed by design and the quartic coefficient is what remains. Now the same design wearing a $1^{\infty}$ costume:
\[
\lim_{x\to0}\left(\frac{\cosh x+\cos x}{2}\right)^{1/x^{4}}=\eu^{1/24}.
\]
The averaging kills the $x^{2}$ terms, since $\cosh$ contributes $+\tfrac{x^{2}}{2}$ and $\cos$ contributes $-\tfrac{x^{2}}{2}$, while the quartic terms reinforce: both contribute $+\tfrac{x^{4}}{24}$, so the base is $1+\tfrac{x^{4}}{24}+O(x^{8})$. The $1^{\infty}$ shortcut gives $\exp\bigl(\lim x^{-4}\cdot\tfrac{x^{4}}{24}\bigr)=\eu^{1/24}$. Verified as $1.0425469$ at $x=0.05$ against $\eu^{1/24}=1.0425469$.
\end{worked}

\begin{trap}
Attacking the first form with four rounds of L'Hopital. It works, but the fourth derivative of the numerator has to be produced without error and the form has to be rechecked at each of the four stages. Reading one coefficient off the table is not merely faster, it is a different order of reliability.
\end{trap}

\begin{seealsob}Expand to the first surviving order; the standard expansion table; the exponential conversion for power forms.\end{seealsob}

\begin{keynote}[Read the denominator as a hint]
In a constructed problem the exponent of the denominator is not arbitrary: it is chosen so that the limit is finite and nonzero, which means it equals the order of the numerator. So the denominator tells you in advance how deep the cancellation goes, and therefore how far to expand. If your expansion produces a numerator of a different order than the denominator, the likeliest explanation is an arithmetic slip, not an infinite or zero limit. Check before believing.
\end{keynote}

The engineered cases are generous: the denominator announces how deep the cancellation goes. The remaining case is the ungenerous one, where the cancellation is deeper than any reasonable person would guess, and the only defence is to expand honestly and keep going.

\tech{When the third order is genuinely needed}{hard}
\cue The denominator is $x^{3}$ or higher, or the numerator is a difference of two functions that agree to second order and possibly much further.
\method Push both expansions to exactly the order the denominator demands, using the budget rule, and do not stop early because several orders have cancelled. Some classics need far more than three: $\sin(\tan x)$ and $\tan(\sin x)$ agree through order $6$, so anything built from their difference lives at order $7$.

\begin{worked}
\[
\lim_{x\to0}\frac{\sin(\tan x)-\tan(\sin x)}{x^{7}}.
\]
Both composites expand to $x+\tfrac{x^{3}}{6}-\tfrac{x^{5}}{40}+O(x^{7})$, and the two agree at orders $1$, $3$ and $5$; the even orders vanish because both functions are odd. The first surviving difference is at order $7$, where the coefficients differ by $-\tfrac1{30}$, so the limit is $-\dfrac1{30}$. The numerical evidence has to be read carefully: the quotient is $-0.033472$ at $x=0.06$ and $-0.033368$ at $x=0.03$, approaching $-0.033333$ quadratically in $x$, exactly as an $O(x^{2})$ relative correction predicts.
\end{worked}

\begin{trap}
Concluding after two orders of cancellation that ``the limit is $0$'' and stopping. Cancellation through order $k$ says that the answer is the coefficient at order $k+1$; it says nothing about that coefficient being zero. The second trap here is numerical: at $x=0.3$ the quotient reads $-0.0370$, which rounds to neither $-1/30$ nor anything else recognisable, and a solver who trusts it will chase a nonexistent closed form.
\end{trap}

The deep agreement in that example is not an accident, and knowing why saves you from expecting it where it does not occur. Both $\sin\circ\tan$ and $\tan\circ\sin$ are odd, so all even coefficients vanish for free; and the order-$3$ and order-$5$ coefficients of a composite $u\circ v$ built from two members of the $x\pm\tfrac{x^{3}}{k}$ family depend only on the sum of the two signature coefficients, which is symmetric under swapping the functions. Order $7$ is the first place where the composition order can be felt. The same structure gives
\[
\lim_{x\to0}\frac{\arcsin(\arctan x)-\arctan(\arcsin x)}{x^{7}}=-\frac1{30},
\]
the same value, verified as $-0.033318$ at $x=0.03$. Whenever you meet a pair of composites of this shape, budget for order $7$ before you start.

\begin{seealsob}Order bookkeeping across a subtraction; composite and nested expansion; repeated L'Hopital with simplification.\end{seealsob}

\subsection{An order-counting drill}

For each quotient, state $d=\ord D$ and the budget for each numerator factor before computing anything. The point is that all of this is decidable by inspection.

\begin{enumerate}[label=(\roman*)]
\item $\dfrac{\sin x-x}{x^{3}}$. Here $d=3$; one factor, budget $3$; answer $-\tfrac16$.
\item $\dfrac{(1-\cos x)\ln(1+x)}{x^{3}}$. Here $d=3$, orders $2$ and $1$ summing to $3$, so each factor needs only its leading term; answer $\tfrac12$.
\item $\dfrac{\tan x-\arctan x}{x^{3}}$. Here $d=3$ and the linear terms cancel, so both must go to order $3$; the signature terms are $+\tfrac13$ and $-\tfrac13$, giving $\tfrac23$.
\item $\dfrac{\eu^{x^{2}}-\cos x}{x^{2}}$. Here $d=2$; the inner argument $x^{2}$ has order $2$, so $\eu^{x^{2}}=1+x^{2}+O(x^{4})$ needs only one term of the outer series; against $\cos x=1-\tfrac{x^{2}}{2}+O(x^{4})$ the answer is $\tfrac32$.
\item $\dfrac{x^{2}(\eu^{x}-1)}{\sin x\,(1-\cos x)}$. Here $d=1+2=3$, the numerator has order $2+1=3$, and both leading coefficients are available at once; the answer is $\tfrac{1}{1\cdot\frac12}=2$.
\item $\dfrac{1-\cos(x^{2})}{x^{2}\sin^{2}x}$. Composition multiplies orders, so the numerator has order $4$ with leading term $\tfrac{x^{4}}{2}$; the denominator has order $2+2=4$ with leading coefficient $1$; the answer is $\tfrac12$. Nobody needs to expand anything beyond one term.
\item $\dfrac{\arctan x-\operatorname{arsinh}x}{x^{3}}$. Both are $x\mp\tfrac{x^{3}}{k}$ types with signature terms $-\tfrac13$ and $-\tfrac16$, so $d=3$ and the answer is $-\tfrac13+\tfrac16=-\tfrac16$. The linear terms cancel; the cubic terms do not.
\end{enumerate}

Not one of these required a decision to be revised halfway through, and that is the whole benefit. A solver who counts orders first never expands twice.

\section{Moving the expansion point, and dividing series}

Everything so far assumed the limit point was $0$. That assumption costs nothing, because a substitution always restores it, but the substitution has to be made explicitly rather than assumed. The three techniques here move the expansion point to a finite $a$, to infinity, and then handle the one algebraic operation that expansions do not support directly, namely division by a series with no constant term.

\begin{keynote}[The three substitutions]
$h=x-a$ moves a finite limit point to the origin. $t=1/x$ moves $\pm\infty$ to the origin. $u=1/(x-a)$ moves a point where the expression blows up to infinity, which is occasionally what a problem wants. In every case the substitution must be applied to the \emph{whole} expression, exponents included, and the direction of approach must be tracked: $x\to a^{+}$ becomes $h\to0^{+}$, and $x\to-\infty$ becomes $t\to0^{-}$, which is where signs are lost.
\end{keynote}

\tech{Expansion at a point other than zero}{medium}
\cue $x\to a$ with $a\neq0$: most often $x\to1$, $x\to\pi/2$, or $x\to\eu$.
\method Substitute $h=x-a$ and expand in $h$:
\[
f(a+h)=f(a)+f'(a)h+\frac{f''(a)}{2}h^{2}+O(h^{3}).
\]
For the standard functions it is usually quicker to use an addition formula than to differentiate: $\sin(a+h)=\sin a\cos h+\cos a\sin h$, then expand $\cos h$ and $\sin h$ from the table. Once the substitution is made, every technique for $x\to0$ applies verbatim.

\begin{worked}
\[
\lim_{x\to1}\left\{\frac{-ax+\sin(x-1)+a}{x+\sin(x-1)-1}\right\}^{(1-x)/(1-\sqrt{x})},\qquad a<1 .
\]
Put $h=x-1$. The numerator is $a(1-x)+\sin(x-1)=-ah+\bigl(h-\tfrac{h^{3}}{6}+O(h^{5})\bigr)=(1-a)h+O(h^{3})$, and the denominator is $h+\sin h=2h+O(h^{3})$, so the base tends to $\tfrac{1-a}{2}$, a positive number. The exponent factors rather than expands: $\dfrac{1-x}{1-\sqrt x}=\dfrac{(1-\sqrt x)(1+\sqrt x)}{1-\sqrt x}=1+\sqrt x\to2$. The limit is therefore $\Bigl(\dfrac{1-a}{2}\Bigr)^{2}$. Verified at $x=1+10^{-6}$: the value is $0.2500000$ for $a=0$ and $4.000001$ for $a=-3$.

A second, shorter instance:
\[
\lim_{x\to\pi/2}\frac{1-\sin x}{(x-\pi/2)^{2}}=\lim_{h\to0}\frac{1-\cos h}{h^{2}}=\frac12 ,
\]
using $\sin(\pi/2+h)=\cos h$. Verified as $0.5000000$ at $h=10^{-4}$.

And one where the addition formula does the work:
\[
\lim_{x\to\pi/4}\frac{\tan x-1}{x-\pi/4}
=\lim_{h\to0}\frac{\dfrac{1+\tan h}{1-\tan h}-1}{h}
=\lim_{h\to0}\frac{2\tan h}{h(1-\tan h)}=2 ,
\]
using $\tan(\pi/4+h)=\dfrac{1+\tan h}{1-\tan h}$ and $\tan h\sim h$. This is of course $\sec^{2}(\pi/4)=2$, the derivative of $\tan$ at $\pi/4$, which is the point: a difference quotient at $a$ is an expansion at $a$ truncated at first order.
\end{worked}

\begin{trap}
Expanding $\sin(x-1)$ ``around $x=0$'' because that is the table you memorised. The expansion point must be the limit point. A related error is expanding the base around $a$ but leaving the exponent in terms of $x$ and then taking limits of the two separately, which is only legal when the base has a limit different from $1$, as it does here.
\end{trap}

\begin{seealsob}Composite and nested expansion; expansion at infinity via $t=1/x$; the exponential conversion for power forms.\end{seealsob}

The limit point $\infty$ is handled the same way, by a substitution that manufactures a small quantity where none was visible.

\tech[Expansion at infinity via reciprocal substitution]{Expansion at infinity via $t=1/x$}{easy}
\cue $x\to\infty$ with no single dominant term, or a difference of two quantities that both grow and nearly cancel.
\method Set $t=1/x$ and expand in $t\to0^{+}$. The workhorse case is
\[
\sqrt{x^{2}+ax+b}=x\sqrt{1+at+bt^{2}}=x\Bigl(1+\frac{a}{2}t+\frac{4b-a^{2}}{8}t^{2}+O(t^{3})\Bigr)
=x+\frac{a}{2}+\frac{4b-a^{2}}{8x}+O(x^{-2}),
\]
valid for $x\to+\infty$; for $x\to-\infty$ pull out $\abs{x}=-x$ instead.

\begin{worked}
\[
\lim_{x\to\infty}\bigl(\sqrt{x^{2}+x}-x\bigr)=\lim_{x\to\infty}\Bigl(x+\tfrac12+O(x^{-1})-x\Bigr)=\frac12 ,
\]
directly from the display with $a=1$, $b=0$. Verified as $0.4999999$ at $x=10^{6}$. A case that needs the next order:
\[
\lim_{x\to\infty}x^{3/2}\bigl(\sqrt{x+1}+\sqrt{x-1}-2\sqrt{x}\bigr).
\]
Write $\sqrt{x\pm1}=\sqrt x\,(1\pm t)^{1/2}$ with $t=1/x$. The linear terms $\pm\tfrac12t$ cancel in the sum, and the quadratic terms reinforce: $(1+t)^{1/2}+(1-t)^{1/2}=2-\tfrac{t^{2}}{4}+O(t^{4})$. So the bracket is $-\tfrac{\sqrt x}{4x^{2}}+O(x^{-7/2})$ and the limit is $-\dfrac14$, verified as $-0.2500000$ at $x=10^{4}$.
\end{worked}

\begin{trap}
Expanding $\sqrt{x^{2}+x}$ as though $x$ were small. The binomial series needs a small perturbation, and manufacturing one is the entire purpose of the substitution. The second trap is dropping the $O(t^{2})$ term in the first example because ``it goes to zero'': it does, but only after multiplication by the $x$ that was pulled out, so its order must be tracked through that multiplication.
\end{trap}

\begin{seealsob}The binomial series for radicals; divide by the dominant term; radicals at infinity.\end{seealsob}

For a difference of exactly two square roots there is a shortcut worth keeping: multiplying by the conjugate turns $\sqrt A-\sqrt B$ into $(A-B)/(\sqrt A+\sqrt B)$, and since $A-B$ is usually a polynomial the whole problem collapses without any series at all. On $\sqrt{x^{2}+x}-x$ this gives $x\big/\bigl(\sqrt{x^{2}+x}+x\bigr)\to\tfrac12$ immediately. The substitution route is the general one, however, because it survives cube roots, three-term differences, and the case where the polynomial difference is itself of high degree, none of which the conjugate handles.

One operation remains unsupported. A Maclaurin series exists only for a function finite at the origin, so a quotient whose denominator vanishes has no series until the vanishing factor is removed. Removing it is mechanical.

\tech{Series division and reciprocal series}{medium}
\cue A quotient of two vanishing series, or the reciprocal of a series with zero constant term, or a difference of two such reciprocals.
\method Factor the leading power out of the denominator and expand the reciprocal of what remains geometrically, using $\dfrac{1}{1+w}=1-w+w^{2}-\cdots$ with $w$ the sum of all the non-constant terms:
\[
\frac{1}{x\bigl(1+c_1x+c_2x^{2}+\cdots\bigr)}=\frac1x\Bigl(1-c_1x+(c_1^{2}-c_2)x^{2}+\cdots\Bigr).
\]
The extracted $1/x$ is a genuine pole, and in a difference of two such reciprocals the poles are what cancel.

\begin{worked}
Taking $\ln(1+x)=x\bigl(1-\tfrac{x}{2}+\tfrac{x^{2}}{3}-\cdots\bigr)$, so $c_1=-\tfrac12$ and $c_2=\tfrac13$:
\[
\frac{1}{\ln(1+x)}=\frac1x\Bigl(1+\frac{x}{2}+\bigl(\tfrac14-\tfrac13\bigr)x^{2}+\cdots\Bigr)=\frac1x+\frac12-\frac{x}{12}+O(x^{2}),
\]
verified numerically: $\bigl(1/\ln(1+x)-1/x-\tfrac12\bigr)/x$ is $-0.083318$ at $x=10^{-4}$ against $-1/12=-0.083333$. The immediate payoff:
\[
\lim_{x\to0}\left(\frac1x-\frac{1}{\eu^{x}-1}\right)
=\lim_{x\to0}\left(\frac1x-\Bigl(\frac1x-\frac12+\frac{x}{12}+O(x^{2})\Bigr)\right)=\frac12 ,
\]
using $\eu^{x}-1=x\bigl(1+\tfrac{x}{2}+\tfrac{x^{2}}{6}+\cdots\bigr)$ and the same reciprocal formula. Verified as $0.4999917$ at $x=10^{-4}$.

The same extraction applied to the cotangent, which is where it is most often needed:
\[
\cot x=\frac{\cos x}{\sin x}=\frac{1-\tfrac{x^{2}}{2}+O(x^{4})}{x\bigl(1-\tfrac{x^{2}}{6}+O(x^{4})\bigr)}
=\frac1x\Bigl(1-\tfrac{x^{2}}{2}\Bigr)\Bigl(1+\tfrac{x^{2}}{6}\Bigr)+O(x^{3})=\frac1x-\frac{x}{3}+O(x^{3}),
\]
so $\displaystyle\lim_{x\to0}\frac{\cot x-1/x}{x}=-\frac13$, verified as $-0.3333334$ at $x=10^{-3}$. Every $\infty-\infty$ limit built from $\cot$, $\csc$, or $1/\ln(1+x)$ is this calculation wearing different clothes.
\end{worked}

\begin{trap}
Trying to ``expand $1/\ln(1+x)$'' as a Maclaurin series. It has a pole at $0$ and no such series exists; the pole must be extracted first. The second trap is forgetting the $c_1^{2}$ term in the reciprocal formula: $1/(1+w)$ is $1-w+w^{2}$, and $w^{2}$ contributes at order $2$ whenever $c_1\neq0$. In the $\ln(1+x)$ computation above, omitting it turns $\tfrac14-\tfrac13=-\tfrac1{12}$ into $-\tfrac13$, and the resulting answer for any $\infty-\infty$ limit built from it is wrong by a factor of four.
\end{trap}

\begin{seealsob}Choosing the profitable conversion; the $\infty-\infty$ conversion; expand to the first surviving order.\end{seealsob}

\section{L'Hopital's rule, in its proper place}

The rule is a genuine theorem and it is occasionally the shortest path. It is also the most over-applied result in elementary analysis, and the reason is that its statement is easy to remember while its hypotheses are easy to forget. This section states it with the hypotheses attached, shows the one shape where it beats expansion, shows how to keep it from exploding, and then catalogues the three ways it fails.

Keep the keynote from the start of this chapter in view while reading it. One round of the rule is one order of truncation of two Taylor series, so the rule can never produce information that expansion could not. What it can do is produce that information without requiring the series to be known, which matters exactly when the functions are not in the table: at infinity, where Maclaurin series do not apply; at a finite point where the functions are arbitrary rather than standard; and in the rare case where a single differentiation happens to simplify an expression rather than complicate it. Those three cases are the legitimate domain of the rule, and they are narrower than most courses suggest.

\tech[LHopital rule with hypotheses]{The rule, with its hypotheses actually checked}{easy}
\cue Substitution gives $\tfrac00$ or $\tfrac{\infty}{\infty}$, and both functions are differentiable near, though not necessarily at, the point.
\method The theorem: if $f,g\to0$ (or both $\to\pm\infty$) as $x\to a$, if $g'\neq0$ on a punctured neighbourhood of $a$, and if $\lim f'/g'$ exists finitely or as $\pm\infty$, then
\[
\lim_{x\to a}\frac{f(x)}{g(x)}=\lim_{x\to a}\frac{f'(x)}{g'(x)}.
\]
Verify the form at \emph{every} application, not only the first. Differentiate numerator and denominator separately; the quotient rule has no role here.

\begin{worked}
\[
\lim_{x\to0}\frac{\eu^{x}-1-x}{x^{2}}.
\]
Form $\tfrac00$; $g'=2x$ is nonzero on a punctured neighbourhood. One round gives $\dfrac{\eu^{x}-1}{2x}$, which is again $\tfrac00$, and the hypothesis holds again. A second round gives $\dfrac{\eu^{x}}{2}\to\dfrac12$. Since that limit exists, both applications are justified in retrospect, which is the correct logical order: the rule concludes backwards from the existence of the derivative limit. Verified as $0.5000167$ at $x=10^{-4}$. Compare the expansion route: $\eu^{x}-1-x=\tfrac{x^{2}}{2}+O(x^{3})$, one line.

A case where the rule genuinely wins, because there is no Maclaurin series to use:
\[
\lim_{x\to\infty}\frac{\ln x}{\sqrt[3]{x}}
=\lim_{x\to\infty}\frac{1/x}{\tfrac13x^{-2/3}}
=\lim_{x\to\infty}\frac{3}{x^{1/3}}=0 .
\]
The form is $\tfrac{\infty}{\infty}$, $g'=\tfrac13x^{-2/3}$ is nonzero throughout, and one round converts the problem into a determinate one. The growth hierarchy gives the same answer even faster, but this is the shape in which a single round of the rule costs nothing and settles everything.
\end{worked}

\begin{trap}
Differentiating a quotient that is no longer indeterminate. After one round the form is often $\tfrac{c}{0}$ or $\tfrac{c}{d}$, and applying the rule there produces nonsense: on $\dfrac{x}{x+1}$ at $x\to0$, the rule would give $1$ instead of $0$. The other classic is differentiating the quotient rather than the two functions separately.
\end{trap}

\begin{seealsob}The point-value chain; repeated L'Hopital with simplification; the rule's failure modes.\end{seealsob}

There is one shape where a single round of the rule is genuinely the best move: a limit at a finite point where the functions are not standard ones with memorised series, so that expansion would mean computing derivatives anyway.

\tech{The point-value chain}{medium}
\cue A limit at a specific finite point $a$ where the numerator's vanishing is not obvious, and where the functions are not in the Maclaurin table.
\method Apply the rule once and evaluate $f'(a)$ and $g'(a)$ honestly as numbers. The interest of these problems is usually that $f'(a)$ turns out to vanish as well, which settles the limit at once, or that $g'(a)$ is a clean nonzero number, in which case the answer is $f'(a)/g'(a)$ and there is nothing further to do.

\begin{worked}
\[
\lim_{x\to\eu}\frac{\eu^{x}-x^{\eu}}{\eu-x}.
\]
Both numerator and denominator vanish at $x=\eu$, since $\eu^{\eu}=\eu^{\eu}$. Differentiating, $f'(x)=\eu^{x}-\eu\,x^{\eu-1}$ and $g'(x)=-1$. At $x=\eu$,
\[
f'(\eu)=\eu^{\eu}-\eu\cdot\eu^{\eu-1}=\eu^{\eu}-\eu^{\eu}=0,
\]
so the limit is $0/(-1)=0$. Verified: the quotient is $-2.8\times10^{-6}$ at $x=\eu+10^{-6}$, which is consistent with a limit of $0$ approached linearly. The whole problem was the observation that $\eu^{x}$ and $x^{\eu}$ are tangent at $x=\eu$.

When one round is not enough, expanding about $a$ is usually better than a second round. For
\[
\lim_{x\to1}\frac{x^{x}-x}{x-1-\ln x},
\]
put $h=x-1$. Then $x\ln x=(1+h)\bigl(h-\tfrac{h^{2}}{2}+O(h^{3})\bigr)=h+\tfrac{h^{2}}{2}+O(h^{3})$, so $x^{x}=\eu^{h+h^{2}/2+\cdots}=1+h+h^{2}+O(h^{3})$ and the numerator is $h^{2}+O(h^{3})$. The denominator is $h-\ln(1+h)=\tfrac{h^{2}}{2}+O(h^{3})$. The limit is $2$, verified as $2.0000224$ at $x=1+10^{-5}$. Two rounds of L'Hopital would require differentiating $x^{x}$ twice.
\end{worked}

\begin{trap}
Mis-differentiating $x^{\eu}$ as $x^{\eu}\ln x$, which is the rule for $\eu^{x}$, or $\eu^{x}$ as $x\,\eu^{x-1}$, which is the rule for $x^{\eu}$. The problem is engineered around exactly this confusion, and the two wrong derivatives give two different wrong answers, neither of which is $0$.
\end{trap}

\begin{seealsob}The rule, with its hypotheses actually checked; expansion at a point other than zero; repeated L'Hopital with simplification.\end{seealsob}

When one round is not enough, the rule's cost grows faster than most solvers expect, because each differentiation of a product multiplies the number of terms. Controlling that growth is a technique in itself.

\tech[Repeated LHopital with simplification]{Repeated application, with simplification between rounds}{medium}
\cue One round of the rule has left you still at $\tfrac00$, and the expression is a product or a quotient of several factors.
\method Before differentiating again, simplify and cancel. Pull out any factor whose limit exists and is finite and nonzero, and evaluate it separately, using $\lim(uv)=\lim u\cdot\lim v$ when both limits exist. Apply trigonometric identities to collapse the derivative you just produced. Each such step keeps the expression from doubling in length.

\begin{worked}
\[
\lim_{x\to0}\frac{\tan x-x}{x^{3}}.
\]
Raw, this takes three rounds and the third derivative of $\tan$. With one identity it takes one round: differentiating gives $\dfrac{\sec^{2}x-1}{3x^{2}}$, and $\sec^{2}x-1=\tan^{2}x$, so the expression is $\dfrac13\Bigl(\dfrac{\tan x}{x}\Bigr)^{2}\to\dfrac13$. Verified as $0.3333335$ at $x=10^{-3}$. The expansion route is shorter still: $\tan x-x=\tfrac{x^{3}}{3}+O(x^{5})$ from the table, so the answer is the signature coefficient of $\tan$ and nothing needs to be differentiated at all.

A second instance, where the useful simplification is factoring rather than an identity:
\[
\lim_{x\to0}\frac{\eu^{x}\sin x-x-x^{2}}{x^{3}}=\frac13 .
\]
Three rounds of the rule on $\eu^{x}\sin x$ generate eight terms by the product rule. Expanding instead, $\eu^{x}\sin x=\bigl(1+x+\tfrac{x^{2}}{2}+\tfrac{x^{3}}{6}\bigr)\bigl(x-\tfrac{x^{3}}{6}\bigr)+O(x^{4})=x+x^{2}+\bigl(\tfrac12-\tfrac16\bigr)x^{3}+O(x^{4})$, so the numerator is $\tfrac{x^{3}}{3}+O(x^{4})$. Verified as $0.3333333$ at $x=10^{-3}$.
\end{worked}

\begin{trap}
Never simplifying. By the third derivative of a product of four factors you have thirty terms and the limit is lost in the bookkeeping. This is precisely the regime where expansion is strictly better, because a Taylor coefficient is a lookup while a third derivative is a computation.
\end{trap}

\begin{seealsob}Expand to the first surviving order; the rule, with its hypotheses actually checked; the rule's failure modes.\end{seealsob}

Finally, the failures. They are not exotic pathologies invented to embarrass students. Two of the three appear in ordinary problems, and the third is the reason the theorem's statement is phrased as a one-way implication.

\tech{The rule's failure modes}{hard}
\cue $\lim f'/g'$ does not exist, or reproduces the quotient you started with, or the hypothesis on $g'$ is violated.
\method Recognise three distinct failures and the correct response to each.
\emph{(a) Cycling.} $\displaystyle\lim_{x\to\infty}\frac{\sqrt{1+x^{2}}}{x}$ differentiates to $\dfrac{x}{\sqrt{1+x^{2}}}$, which differentiates back to $\dfrac{\sqrt{1+x^{2}}}{x}$: an infinite loop. Response: divide by the dominant term, giving $\sqrt{x^{-2}+1}\to1$.
\emph{(b) The derivative limit does not exist although the original does.} Response: the rule is silent, not negative; use squeeze or dominance instead.
\emph{(c) The hypothesis on $g'$ is violated}, because $g'$ vanishes on a sequence approaching the point. Response: the rule does not apply at all, and no conclusion of any kind may be drawn from $f'/g'$.

\begin{worked}
The standard instance of (b):
\[
\lim_{x\to\infty}\frac{x+\sin x}{x+\cos x}=1,
\]
by dividing by $x$: the quotient is $\dfrac{1+x^{-1}\sin x}{1+x^{-1}\cos x}$ and both perturbations are squeezed to $0$. Verified as $1.0000000130$ at $x=10^{8}$. The derivative quotient, however, is $\dfrac{1+\cos x}{1-\sin x}$, which oscillates forever and, worse, has zero denominator at every $x=\tfrac{\pi}{2}+2k\pi$, so hypothesis (c) fails too.

An instance at a finite point, combining (b) and (c):
\[
\lim_{x\to0}\frac{x^{2}\sin(1/x)}{\sin x}=0,
\]
since $\abs{x^{2}\sin(1/x)}\leq x^{2}$ and $\sin x\sim x$, so the quotient is squeezed between $\pm\abs{x}$. The derivative quotient is $\dfrac{2x\sin(1/x)-\cos(1/x)}{\cos x}$, whose numerator oscillates between $-1$ and $1$ without settling: it reads $-0.872$ at $x=10^{-2}$ and $-0.561$ at $x=10^{-3}$. The rule reports nothing; the limit is $0$ regardless.

Failure (c) in isolation, and it is the sharpest of the three. Take $f(x)=x+\sin x\cos x$ and $g(x)=\eu^{\sin x}\bigl(x+\sin x\cos x\bigr)$ as $x\to\infty$. Both tend to $\infty$, so the form is right. But $f/g=\eu^{-\sin x}$, which oscillates forever between $\eu^{-1}$ and $\eu$ and has no limit. Meanwhile $f'=1+\cos2x=2\cos^{2}x$ and $g'=\eu^{\sin x}\cos x\bigl(x+\sin x\cos x+2\cos x\bigr)$, so
\[
\frac{f'}{g'}=\frac{2\cos x}{\eu^{\sin x}\bigl(x+\sin x\cos x+2\cos x\bigr)}\longrightarrow 0 .
\]
The derivative quotient converges beautifully to $0$ and the original quotient has no limit at all. The rule is not violated, because $g'$ vanishes at every $x=\tfrac{\pi}{2}+k\pi$, so the hypothesis was never satisfied. Numerically: $f/g$ reads $1.659$ at $x=100$ and $0.437$ at $x=1000$, while $f'/g'$ reads $0.0283$ and $0.00049$ at the same points. The moral is that ``$g'\neq0$ on a punctured neighbourhood'' is a condition to be checked, not a formality, and that a well-behaved derivative quotient is no evidence whatever that the checking can be skipped.
\end{worked}

\begin{trap}
Concluding from ``$\lim f'/g'$ does not exist'' that ``$\lim f/g$ does not exist''. The theorem is a one-way implication and its converse is false, as the two examples above show. The practical rule: if the derivative quotient is not visibly better than what you started with, abandon the rule immediately rather than differentiating again in hope.
\end{trap}

\begin{seealsob}Divide by the dominant term; the squeeze theorem; repeated L'Hopital with simplification.\end{seealsob}

\begin{keynote}[The standing order]
Expand first. Use L'Hopital when the functions are not in the table and a single differentiation produces a number, as in the point-value chain. Never apply it more than twice without asking what the corresponding Taylor coefficient would have been.
\end{keynote}

The chapter has one thesis and one procedure. The thesis is that a limit at a point is a question about the first few Taylor coefficients, so the method that produces those coefficients directly is the right method, and every other method is either a way of computing the same coefficients more slowly or a way of computing fewer of them. The procedure is: substitute so that the limit point is $0$; compute the order of the denominator; budget each factor of the numerator against that order; expand exactly that far with an explicit remainder; read off the coefficient. L'Hopital's rule enters only when the functions are not in the table, and it leaves as soon as the derivative quotient stops improving.

\section{Recognition matrix}
\begin{recog}{@{}p{0.30\textwidth}p{0.36\textwidth}p{0.28\textwidth}@{}}
\rowcolor{mist}\mh{If you see} & \mh{Reach for} & \mh{If that fails} \\
\midrule
\endhead
$\tfrac00$ at $x=0$ with standard functions & The Maclaurin table, expanded to $\ord$ of the denominator & Push one order further and recheck \\
A radical of something near $1$ & $(1+u)^{a}$ with $u$ the whole perturbation & Second order; radicals over $x^{2}$ need it \\
A function of a function & Expand the inner to $u$, substitute into the outer & Factor out a common exponential first \\
A product in the numerator & Budget: expand $f_i$ through degree $d-(m-m_i)$ & Recompute $m_i$ after any cancellation \\
Two functions with the same leading term, subtracted & Expand both one order past the expected cancellation & Confirm the surviving coefficient is nonzero \\
$f(x)$ minus its own first $m$ Taylor terms & Read off the next coefficient of $f$ & Check whether the intervening one vanishes \\
Cancellation past order $3$, or odd composites & Expand to order $7$; do not conclude $0$ early & Read numerics as a rate, not a value \\
$x\to a$ with $a\neq0$ & $h=x-a$, then an addition formula & Differentiate at $a$ directly \\
$x\to\infty$ with near-cancellation & $t=1/x$, then the binomial series & Conjugate multiplication \\
A reciprocal of a vanishing series & Extract the pole, then expand $1/(1+w)$ & Common denominator instead \\
$\tfrac00$ with non-tabulated functions & One L'Hopital round, then evaluate $f'(a)$, $g'(a)$ & Expand $f$ and $g$ about $a$ by hand \\
The derivative quotient repeats itself & Stop; divide by the dominant term & Substitute $t=1/x$ \\
The derivative quotient oscillates & Stop; the rule is silent, use squeeze & Bound above and below and compare \\
$g'$ vanishes arbitrarily near the point & The rule does not apply; draw no conclusion & Direct estimate or squeeze \\
\end{recog}
